\begin{abstract}
Knowledge of an object's center of mass (CoM) is essential for multi-robot pushing manipulation, yet it is commonly assumed to be known in advance. This paper proposes a distributed method for estimating the CoM of a planar polygonal object using a team of omnidirectional mobile robots that interact with the object solely through pushing. The robots are coordinated to exert zero net force and a nonzero torque on the object, inducing approximately pure rotation about its CoM, which then coincides with the instantaneous center of rotation. This yields a linear measurement model relating each robot's contact point position and velocity and the object's angular velocity to the unknown CoM. The problem is decomposed into formation synthesis and force control. Formation synthesis seeks a robot configuration on the object boundary that maximizes the robustness of the achievable torque range under zero net force, subject to force bounds and collision avoidance, and is solved by augmented Lagrangian particle swarm optimization. In force control, each robot solves a bounded quadratic program to determine its pushing force. During the induced rotation, each robot independently estimates the CoM via recursive least squares without sharing measurements. Simulations with six robots pushing a U-shaped object under measurement noise show that the estimation errors of all robots fall below 1 cm within approximately 0.6 s and settle at around 1 mm, with all robots converging to nearly the same estimate.
\end{abstract}

\begin{IEEEkeywords}
Center of mass estimation, cooperative manipulation, formation synthesis, multi-robot systems, particle swarm optimization, pushing, recursive least squares.
\end{IEEEkeywords}
